% Section 6: Application: electricity dispatch
\subsection{Application: electricity dispatch with decision-focused forecasting}
\label{sec:app_electricity}

A microgrid operator must decide how to schedule the next $24$ hours of
electricity generation, energy-storage charging and discharging, and
demand-side flexibility under uncertainty about renewable generation,
load demand, and price. Each scheduling decision $a$ is a vector of
hour-by-hour set-points subject to ramping constraints, state-of-charge
limits, and exchange caps with the main grid. The realised operating
cost depends on which generators were committed and how the realised
renewable output and load tracked the forecast. The downstream problem is
a two-stage robust optimisation. The forecast object is the joint
distribution of renewable output and load demand at $15$-minute
resolution. The application is the cleanest empirical
decision-focused win in the literature because the downstream
optimisation has a clean linear-plus-convex structure that the SPO
machinery of Section~\ref{sec:dfl} fits exactly.

\citet{CaoXu2025} embed a multilayer probabilistic encoder-decoder
forecaster inside a two-stage robust optimisation solved by the
column-and-constraint generation algorithm. The architecture is a
hierarchy of residual blocks followed by parallel encoder-decoder chains
with a temporal-attention decoder. The training loss combines the
continuous ranked probability score with a regret-type decision-focused
term,
\begin{equation}
\mathrm{Regret}_{\mathrm{SPO}}(y, \hat y)
= J\bigl(\bar z_\rho(\hat y), y\bigr) - J\bigl(z^\star(y), y\bigr),
\label{eq:caoxu_regret}
\end{equation}
where $J(z, y)$ is the realised operating cost of dispatch $z$ under
realisation $y$, $z^\star(y)$ is the oracle dispatch, and $\bar
z_\rho(\hat y)$ is the dispatch produced by a Tikhonov-regularised convex
relaxation of the two-stage robust programme.\footnote{The downstream
problem has binary unit-commitment variables and is therefore
non-differentiable. \citet{CaoXu2025} relax the binaries to $[0, 1]$ and
add $\rho \|z\|^2$ to guarantee strong convexity. Gradients of
\eqref{eq:caoxu_regret} are obtained by implicit differentiation of the
Karush-Kuhn-Tucker residual through the relaxed problem, in the spirit
of \citet{DontiAmosKolter2017}. The forward decision evaluation still uses
the exact mixed-integer programme via column-and-constraint generation,
and only the gradient path runs through the convex surrogate.}

The empirical comparison is on the modified IEEE 33-bus and 69-bus power
distribution networks, each with two energy-storage units and several
photovoltaic and wind generators. Table~\ref{tab:caoxu_breakdown} shows
the daily cost breakdown on the 33-bus system across four
forecasting-optimisation schemes. The deterministic-forecast pipeline
costs \$5{,}127 per day. Replacing the deterministic forecast with a
DeepAR probabilistic forecaster shaves $5.9\%$. Switching to the
\citet{CaoXu2025} encoder-decoder (MED) trained on CRPS shaves
$11.5\%$. Adding the decision-focused regret loss to the MED training
brings the reduction to $18.0\%$.\footnote{On the larger 69-bus system the
corresponding net-cost reductions are $4.6\%$, $8.4\%$, and $12.6\%$ from
a baseline of \$7{,}521 per day. The forecasting-only improvements (DeepAR
and MED without SPO) account for roughly two thirds of the total
reduction; the decision-focused fine-tuning contributes the remaining
third.}

\begin{table}[h]
\centering
\caption{Daily net operating cost (\$) on the IEEE 33-bus microgrid for
four forecasting-optimisation schemes \citep{CaoXu2025}. The
decision-focused MED-plus-SPO training delivers an $18.0\%$ reduction over
the deterministic baseline.}
\label{tab:caoxu_breakdown}
\begin{tabular}{lcc}
\hline\hline
Scheme              & Net daily cost (\$) & Reduction vs baseline \\
\hline
Deterministic + TSRO      & 5{,}127            & --                  \\
DeepAR + TSRO             & 4{,}823            & $5.9\%$             \\
MED + TSRO (CRPS only)    & 4{,}536            & $11.5\%$            \\
MED + SPO + TSRO          & 4{,}205            & $18.0\%$            \\
\hline\hline
\end{tabular}
\end{table}

The forecast-side gains and the decision-side gains decompose cleanly.
The MED architecture's CRPS reduction over DeepAR captures
forecast-accuracy improvements that the operator could in principle have
obtained without changing the optimisation routine. The MED-plus-SPO
fine-tuning captures a further decision-aware adjustment that does not
improve the CRPS but does improve realised operating cost.
\citet{CaoXu2025} also report a $42\%$ narrowing of the $95\%$
confidence interval on total operation cost under a stress-test of the
uncertainty budget $\Gamma$, the calibration improvement that the SPO
programme of Section~\ref{sec:dfl} is theoretically designed to deliver.
The chapter's thesis statement carries over directly. The proposed method
aligns forecasting accuracy with operational objectives by directly
minimising decision regret via a surrogate SPO loss.
\citet{DontiAmosKolter2017} on PJM load data, mentioned in
Section~\ref{sec:dfl}, is the antecedent at the larger-scale grid level
and reaches the same $38.6\%$-versus-baseline qualitative conclusion.

The microgrid case study is the chapter's cleanest piece of evidence that
decision-focused training is not a methodological curiosity but an
operational lever for cost reduction in problems with both a quantifiable
forecast object and a structured downstream decision.
